\section{System Architecture}

\subsection{Layered View}
RETROSPECT is organized as a three-layer Cloud--Fog--Edge system in which each layer has a clearly delimited operational responsibility.
\begin{itemize}
    \item \textbf{Cloud layer.} This layer hosts the API server, dashboard, controllers, CRDs, and supporting services on Kubernetes. It is responsible for storing desired state, exposing operator-facing APIs, reconciling resources, and persisting reported execution state.
    \item \textbf{Fog layer.} This layer hosts the gateway and the emulation support services. It acts as the operational bridge between cluster-native control logic and embedded endpoints, translating orchestration events into southbound protocol messages and vice versa.
    \item \textbf{Edge layer.} This layer includes Zephyr-based firmware, the WAMR execution environment, and the deployed WASM application artifacts. It is the execution domain in which the application's intended behavior materializes.
\end{itemize}

The layered decomposition is important not only conceptually but also operationally. It localizes complexity where it is most affordable. Cloud components retain the full expressive power of Kubernetes, the gateway absorbs protocol and coordination complexity, and the device remains lightweight enough to match embedded constraints.

\begin{figure*}[t]
\centering
\begin{tikzpicture}[
  node distance=0.8cm and 0.8cm,
  box/.style={draw,rounded corners,align=center,minimum width=2.6cm,minimum height=0.9cm,font=\footnotesize},
  arr/.style={-Latex,thick}
]

\node[box,fill=blue!8] (dash) {Dashboard\\(Web UI)};
\node[box,fill=blue!8,right=of dash] (api) {API Server\\(Rust)};
\node[box,fill=blue!8,right=of api] (ctrl) {Controllers\\Device/App/Gateway};
\node[box,fill=blue!8,right=of ctrl] (crd) {CRDs\\Device/Application/Gateway};

\node[box,fill=green!10,below=1.3cm of api] (gw) {Gateway\\TLS + CBOR Hub};
\node[box,fill=green!10,left=of gw] (ren) {Renode Manager};

\node[box,fill=orange!10,below=1.3cm of gw] (zep) {Zephyr Firmware};
\node[box,fill=orange!10,right=of zep] (wamr) {WAMR Runtime};
\node[box,fill=orange!10,right=of wamr] (app) {WASM App};

\draw[arr] (dash) -- (api);
\draw[arr] (api) -- (ctrl);
\draw[arr] (ctrl) -- (crd);
\draw[arr] (api) -- (ren);
\draw[arr] (ren) -- (zep);
\draw[arr] (zep) -- (wamr);
\draw[arr] (wamr) -- (app);
\draw[arr] (zep) -- node[right,font=\scriptsize]{TLS + CBOR} (gw);
\draw[arr] (gw) -- node[left,font=\scriptsize]{status/update} (crd);

\node[draw,dashed,fit=(dash)(api)(ctrl)(crd),inner sep=6pt,label=above:{\small Cloud Layer}] {};
\node[draw,dashed,fit=(ren)(gw),inner sep=6pt,label=left:{\small Fog Layer}] {};
\node[draw,dashed,fit=(zep)(wamr)(app),inner sep=6pt,label=left:{\small Edge Layer}] {};

\end{tikzpicture}
\caption{RETROSPECT high-level architecture.}
\label{fig:architecture}
\end{figure*}

\subsection{Gateway-Centric Control}
A defining architectural choice is gateway-centric mediation. Devices do not directly manipulate Kubernetes resources, and they do not need to embed Kubernetes-specific logic. Instead, the gateway terminates southbound TLS sessions, validates enrollment messages, associates transport sessions with device identities, dispatches application life-cycle commands, and reports back execution and liveness information to the cloud side.

This mediation has several benefits. First, it prevents constrained devices from carrying the complexity and attack surface of direct control-plane integration. Second, it creates a single point at which protocol policy, identity checks, and message normalization can be enforced. Third, it gives the cloud side a stable integration surface even if the device protocol evolves. The cost is an additional dependency on gateway correctness and availability, which is explicitly discussed later as an engineering trade-off.

\subsection{Resource Model}
RETROSPECT relies on three primary CRDs.
\begin{itemize}
    \item \textbf{Device.} Encodes device identity, gateway preference, runtime attributes, expected public key, and reported connectivity state.
    \item \textbf{Application.} Encodes desired deployment intent for a WASM artifact together with status fields that reflect execution outcomes at the device side.
    \item \textbf{Gateway.} Encodes endpoint metadata, capability descriptors, and the operational information required for devices and controllers to resolve the proper mediation path.
\end{itemize}

The resource model separates desired state from reported state. This is a critical design choice because it allows the system to represent intent even when devices are offline, temporarily unreachable, or in transition. It also makes failure analysis easier, since operators can compare what the platform wants with what the platform has actually observed.

\subsection{Control and Observability Boundaries}
The architecture deliberately distinguishes control ownership from evidence ownership. Operators and higher-level services own desired state changes through Kubernetes resources. The gateway owns protocol-level observations and transforms them into platform-level status updates. Controllers own reconciliation and consistency enforcement. This separation avoids ambiguous responsibility and reduces the likelihood that multiple components will race to write conflicting status information. In practice, the need for these boundaries became especially clear during validation, when false disconnection behavior emerged from an overly aggressive interpretation of incomplete connectivity evidence.
